\chapter{Conclusion}

\section{Bilan du projet}

Ce travail de Bachelor visait à concevoir et réaliser un prototype capable
de mesurer la réserve de marche d'un mouvement horloger de manière
accélérée, tout en explorant un mode de vieillissement complémentaire. Le
prototype réalisé remplit cet objectif~: le remontage, le dévidage contrôlé
du barillet, la détection optique de l'arrêt et le calcul de la réserve de
marche fonctionnent de manière cohérente, et ont été validés
expérimentalement sur le mouvement de référence.

La grande majorité des exigences du cahier des charges sont satisfaites,
souvent avec une marge confortable~: précision de la mesure, protection du
mouvement, sécurité de l'opérateur, autonomie, encombrement. Quelques points
reposent encore sur des calculs ou des extrapolations plutôt que sur des
mesures directes, comme le couple moteur ou la reproductibilité de la mesure.
Néanmoins, rien ne laisse penser à un problème sur ces aspects.

Deux points n'ont en revanche pas pu être menés à terme comme prévu au
départ. L'accouplement automatique, censé éviter d'avoir à engager la
fourche à la main, n'a pas fonctionné et a dû être abandonné au profit d'un
accouplement manuel. Le pilotage simultané de plusieurs postes n'a pas non
plus pu être testé~: seul un module physique a été réalisé et le logiciel
ne permet pas encore de coordonner plusieurs postes en parallèle.

Le passage de ce prototype à une machine de production à 10 postes demande
encore plusieurs évolutions~: une structure commune aux 10 modules, la
gestion du câblage et des alimentations, un bus USB pour les caméras et
l'adaptation du logiciel qui va avec. Le système de fixation du mouvement
devra également être revu pour intégrer la boîte du mouvement plutôt que de
le serrer directement. En effet, il s'agit d'une piste identifiée trop tard pour être intégrée à ce prototype, mais facilement réalisable sur une machine conçue avec cette
solution dès le départ.

\section{Regard personnel}

Ce travail de fin de formation s'est montré plus qu'intéressant. Mettre en
pratique, dans le cadre d'un projet concret, les compétences acquises durant
ces trois années est à mes yeux une belle manière de clôturer mon Bachelor.
J'ai particulièrement apprécié la pluridisciplinarité de ce travail~:
électronique, informatique, mécanique.
D'un point de vue personnel, je suis très satisfait du travail que j'ai pu
réaliser durant le temps mis à disposition.

Ce travail a également constitué mon premier contact avec le milieu
professionnel, n'ayant jamais travaillé auparavant. Je pense qu'il m'a déjà
permis d'apprendre et d'évoluer, notamment en ce qui concerne la gestion de
projet.


% La signature est optionnelle et pas forcément nécessaire...
\vfil
\hspace{8cm}\makeatletter\@author\makeatother\par
\hspace{8cm}\begin{minipage}{5cm}
    % Place pour signature numérique
    \printsignature
\end{minipage}

